% ----------------------------------------------------------
% ELEMENTOS PRÉ-TEXTUAIS
% ----------------------------------------------------------
% \pretextual

% ---
% Capa
% ---
\imprimircapa
% ---

% ---
% Folha de rosto
% (o * indica que haverá a ficha bibliográfica)
% ---
\imprimirfolhaderosto
% ---

% ---
% Inserir folha de aprovação
% ---

% Isto é um exemplo de Folha de aprovação, elemento obrigatório da NBR
% 14724/2011 (seção 4.2.1.3). Você pode utilizar este modelo até a aprovação
% do trabalho. Após isso, substitua todo o conteúdo deste arquivo por uma
% imagem da página assinada pela banca com o comando abaixo:
%
% \includepdf{folhadeaprovacao_final.pdf}
%
\begin{folhadeaprovacao}

  		\includegraphics[width=1\textwidth]{imagens/unitins.png}
  		
  		\ABNTEXchapterfont\large   1. INFORMAÇÕES DO ACADÊMICO:
				
  		\normalsize Nome: Matheus José Alves Silva Santos
  		\tab Matrícula:
  
  		Período 7	
  		
  		E-mail: matheusetf@gmail.com
  		\tab Telefone: (63) 99285-7729
  		
  		\par
  		\vspace*{0.5cm}     
  		  
  		\ABNTEXchapterfont\large   2. INFORMAÇÕES DO ESTÁGIO:
		  		\normalsize 
		
		Professor Orientador: Douglas Chagas da Silva
		
		Início do Estágio: \tab \tab Término do Estágio: 16/06/2018
		
		Total de horas semanais dedicadas ao estágio supervisionado: 20 horas
		
		Área de realização do estágio: Infraestrutura Ágil
		
		
		Data: \_\_\_\_$/$\_\_\_\_$/$\_\_\_\_
		\par
		\vspace*{0.5cm}

  		\ABNTEXchapterfont\large   3. PARECER DO PROFESSOR ORIENTADOR:
  		\normalsize 
  		
  		Aprovado $($ \tab $)$ \tab \tab Reprovado $($ \tab $)$  \tab \tab Nota: \_\_\_\_\_
  		
  	 OBSERVAÇÕES:
  		  		
  		  		 DATA: \_\_\_\_$/$\_\_\_\_$/$\_\_\_\_ \tab \assinatura{Professor Orientador}
  		\par
  		\vspace*{0.5cm}
  		 \ABNTEXchapterfont\large   4. PARECER DO COORDENADOR DE ESTÁGIO:
  		   		\normalsize 
  		
  		 OBSERVAÇÕES: 

  		 DATA: \_\_\_\_$/$\_\_\_\_$/$\_\_\_\_ \tab \assinatura{Coord. de Estágio Supervisionado}
  		 
   \begin{center}
    \vspace*{0.1cm}
    {\large\imprimirlocal}
    \par
    {\large\imprimirdata}
    \vspace*{0.5cm}
  \end{center}
  
\end{folhadeaprovacao}
% ---

% ---
% Dedicatória
% ---
\begin{dedicatoria}
   \vspace*{\fill}
   \centering
   \noindent
   \textit{ Este trabalho é dedicado à minha família, pelo apoio incondicional.} \vspace*{\fill}
\end{dedicatoria}
% ---

% ---
% Agradecimentos
% --- Es
\begin{agradecimentos}
À Deus, pela graça de permitir o meu crescimento no campo de estudo em que me sinto feliz e realizado. \\

Aos meus pais e meus irmãos que sempre apostaram em minha capacidade de crescimento e superar as barreiras que a vida impõe.\\

À minha esposa que tem sido compreensiva nos momentos de dificuldades, e companheira nos desafios que surgem.\\

Aos meus amigos e colegas de universidade, que sempre contribuíram para o desenvolvimento do meu conhecimento e pelas parcerias nos projetos e estudos.


\end{agradecimentos}
% ---

% ---
% Epígrafe
% ---
\begin{epigrafe}
    \vspace*{\fill}
	\begin{flushright}
		\textit{``Não só isso, mas também nos gloriamos \\
			nas tribulações, porque sabemos que a tribulação\\
			produz perseverança; a perseverança, um caráter \\
			aprovado; e o caráter aprovado, esperança. \\
			E a esperança não nos decepciona, porque Deus \\
			derramou seu amor em nossos corações, por meio \\
			do Espírito Santo que Ele nos concedeu.``\\
			(Bíblia Sagrada, Romanos 5:3-5)}
			
	\end{flushright}
\end{epigrafe}
% ---

% ---
% RESUMOS
% ---

% resumo em português
\setlength{\absparsep}{18pt} % ajusta o espaçamento dos parágrafos do resumo
\begin{resumo}
Este projeto visa apresentar um estudo sobre as diversas ferramentas de automação da entrega de software, visando estabelecer uma proposta de cenário que viabilize esse processo dentro de uma cultura DevOps. O cenário deve funcionar como um provedor de serviços sobre plataforma ágil.
 

 \textbf{Palavras-chaves}: Automação, DevOps, Modelo Ágil.
 
\end{resumo}

% resumo em inglês



% ---
% inserir lista de ilustrações
% ---
\pdfbookmark[0]{\listfigurename}{lof}
\listoffigures*
% ---

% ---
% inserir lista de tabelas
% ---
% \newpage
% \pdfbookmark[0]{\listtablename}{lot}
% \newpage
% \listoftables*
% \cleardoublepage
% ---

% ---
% inserir lista de abreviaturas e siglas
% ---
\begin{siglas}
  \item DevOps - Development and Operational (Desenvolvimento e Operacional).
  \item AWS - Amazon Web Services
  \item SSH - Secure Shell
  \item NETCONF - Network Configuration
\end{siglas}
% ---


% ---
% inserir o sumario
% ---
\pdfbookmark[0]{\contentsname}{toc}
\tableofcontents*
\cleardoublepage
% ---

